\section{A local upper bound for the actual prolate defect}

Throughout this section \(\chi\) is a fixed nonprincipal primitive
Dirichlet character, \(R=\Lambda/q\), \(c=\Lambda R\), and
\(\Pi_\Lambda\) is the original full selected raw prolate projection.
The target projection \(Q_\Lambda\), the sharp log window \(S_\Lambda\),
and \(D_\Lambda=S_\Lambda-S_\Lambda Q_\Lambda S_\Lambda\)
are those of AC-T30--T31.
Raw Mellin coordinates are given by the unitary
\((U\psi)(t)=e^{t/2}\psi(e^t)\).
The Sobolev norm \(H_D^s\) means the ordinary \(H^s\) norm
after this unitary.  For a symbol \(a\), \(C_a\) denotes the
logarithmic Fourier multiplier with symbol \(a(\tau)\), using
the positive-sign Fourier convention of AC-T31.

\begin{lemma}[AC-T38: graph approximation with a fixed cutoff margin]
Choose an even \(p\in C_c^\infty((-1,1))\), with \(0\leq p\leq1\),
which equals one on a neighborhood of zero.  Put
\[
 P_\Lambda^p=p(u/\Lambda),\qquad
 Q_R^p=\mathcal F_\epsilon^{-1}p(\xi/R)\mathcal F_\epsilon.
 \tag{TU.1}
\]
Let \(T_\Lambda\) be a Hilbert--Schmidt operator into the raw
parity space with
\(\|T_\Lambda\|_2=O((1+\log\Lambda)^d)\) for some fixed \(d\).
Then, for every \(A>0\),
\[
 \left\|E_\chi(1-\Pi_\Lambda)
             P_\Lambda^pQ_R^pP_\Lambda^pT_\Lambda\right\|_2
                    =O_A(c^{-A}).
 \tag{TU.2}
\]
No form of RH is needed for this lemma.
The same graph approximation, at every polynomial rate, holds
for any fixed finite family of smooth parity profiles compactly
supported strictly inside the raw time interval, after dilation
to radius \(\Lambda\), or compactly supported strictly inside
the raw Fourier interval, after dilation to radius \(R\).
\end{lemma}
\begin{proof}
Let \(b<1\) contain the support of \(p\).
The Fourier kernel of multiplication by \(p(u/\Lambda)\) is
\(\Lambda\widehat p(\Lambda(\xi-\eta))\).
For \(|\eta|\leq bR\) and \(|\xi|>R\), their separation
is at least \((1-b)R\).
Schur's test and rapid decay of \(\widehat p\) consequently give
\[
 \|1_{|\xi|>R}\mathcal F_\epsilon
                         P_\Lambda^pQ_R^p\|
                  \leq C_M c^{-M}\quad(M>0).
 \tag{TU.3}
\]
The constants depend on the fixed \(p,M\), not on the cutoffs.
The full-cutoff prolate inequality (TB.82) applies to the
columns \(f\) of \(P_\Lambda^pQ_R^pP_\Lambda^pT_\Lambda\);
their raw time tail is zero.
Thus their summed squared projection error is bounded by
eight times the squared Fourier-tail norm in (TU.3).
It follows that the raw Hilbert--Schmidt error is
\(O_M(c^{-M}(1+\log\Lambda)^d)\).

For \(0<u<1/2\), the sharp-endpoint Sobolev estimate
(TB.33)--(TB.34), with the full retained eigenvalues bounded
below, gives
\[
 \|\Pi_\Lambda f\|_{H_D^u}
                  \leq C_u(1+c)^u\|f\|_2.
 \tag{TU.4}
\]
For the smooth-cutoff vector itself, differentiate
\(u^{1/2}p(u/\Lambda)(Q_R^pP_\Lambda^p f)(u)\).
The term from \(p'\) is bounded because
the raw support has \(u/\Lambda\) in a fixed compact set;
the band derivative is bounded by \(2\pi R\), multiplied
by at most \(\Lambda\).
Hence its \(H_D^1\) norm is at most \(C(1+c)\|f\|_2\).
Interpolation gives the same estimate (TU.4) for this
smooth-cutoff vector.  Applying these inequalities to an
orthonormal input basis gives their Hilbert--Schmidt versions.

Use the unconditional multiplier bound already used in
AC-T25,
\[
 \|E_\chi f\|_2\leq C_\chi\|f\|_{H_D^{5/16}},
 \tag{TU.5}
\]
and interpolate the small raw \(L^2\) error with the coarse
\(H_D^{3/8}\) bound.  Since
\((5/16)/(3/8)=5/6\), the resulting bound is
\(O_M(c^{-M/6+5/16}(1+\log\Lambda)^d)\).
Choosing \(M\) arbitrarily large proves (TU.2).
For a fixed smooth compact time profile, its Fourier tail
has every polynomial decay rate, while (TU.4) remains
available.  This proves the first profile assertion.
For a compact smooth Fourier profile the time tail has
every such rate and its Fourier tail is zero.
Fourier conjugation changes the dilation generator to its
negative, so its fixed-profile \(H_D^{3/8}\) norm is bounded
after the stated unitary scaling.  The same interpolation
proves the second assertion.
\end{proof}

\begin{theorem}[AC-T39: local convergence of the actual positive defect]
Assume \(\mathsf A\).  For every \(a\in C_c^\infty(\mathbb R)\),
\[
 \lim_{\Lambda\to\infty}
    \operatorname{Tr}\bigl(D_\Lambda C_aC_a^*\bigr)
                =\sum_\gamma m_{\chi,\gamma}|a(\gamma)|^2.
 \tag{TU.6}
\]
Equivalently, the positive measures \(\mu_\Lambda\)
defined in AC-T31 converge vaguely to
\(\sum_\gamma m_{\chi,\gamma}\delta_\gamma\).
This statement concerns compact Mellin-frequency tests.
It does not assert the uniform high-frequency tightness
needed to pass to every original compact norm test.
\end{theorem}
\begin{proof}
Write \(L=\log\Lambda\), \(S_L=1_{[-L,L]}\).
Choose a fixed compact interval containing the support of
\(a\) in its interior, whose endpoints are not zero ordinates,
and list the finitely many distinct ordinates in that interval.
Let \(Z_L\) be the span in \(S_LL^2(\mathbb R)\) of
\[
       1_{[-L,L]}(t)e^{-i\gamma t}t^j,
                         \qquad 0\leq j<m_{\chi,\gamma},
 \tag{TU.7}
\]
and let \(P_{Z,L}\) be its orthogonal projection.
Set \(W_L=S_L-P_{Z,L}\).
The normalized-window Gram calculation in AC-T31 shows that
its dimension is constant for large \(L\), that its normalized
basis coefficients are uniformly bounded, and that
\[
       \|S_LC_aP_{Z,L}\|_2^2
              \longrightarrow
                  \sum_\gamma m_{\chi,\gamma}|a(\gamma)|^2.
 \tag{TU.8}
\]
Indeed each normalized exponential-polynomial jet has Fourier
mass concentrating at its ordinate, so multiplication by the
bounded continuous \(a\) is asymptotically multiplication by
\(a(\gamma)\).  The distinct-ordinate Gram blocks tend to zero,
and each same-ordinate block is the invertible Hilbert matrix
already computed in AC-T31.

\emph{A finite-pole inverse with both exact boundary descriptions.}
On the compact frequency interval, \(a(\tau)/
L(1/2+i\tau,\chi)\) has only the finitely many zero poles
in (TU.7).  Give each pole the boundary prescription
\((\tau-\gamma-i0)^{-j}\).
This defines a compactly supported distribution \(b\).
Its inverse Fourier transform \(\kappa\) is smooth and has
the following two properties, with all derivatives:
\[
 \begin{cases}
 \kappa(t)\text{ decreases faster than every power as }t\to+\infty,\\
 \kappa(t)=\displaystyle\sum_{\gamma,j}d_{\gamma j}
                  t^je^{-i\gamma t}+O_N((1+|t|)^{-N}),
                                      &t\to-\infty .
 \end{cases}
 \tag{TU.9}
\]
Here the polynomial degrees are less than the corresponding
zero multiplicities.  To verify (TU.9), subtract the finitely
many principal parts using smooth cutoffs supported near each
pole.  The remainder is smooth compactly supported.
The inverse transform of each prescribed pole is a constant
multiple of \(1_{t<0}t^{j-1}e^{-i\gamma t}\).
Convolution with the Schwartz transform of its smooth cutoff
gives its full polynomial asymptote on the negative half-line
and a Schwartz remainder at both ends.  This also proves the
assertions for every derivative.

Define
\[
                T_L=U^{-1}C_bW_L.
 \tag{TU.10}
\]
This is an ordinary Hilbert--Schmidt operator, despite the
distribution notation for \(b\).
The exact zero-jet conditions on \(W_Lf\) cancel every
polynomial term of (TU.9) when the output coordinate is
below \(-L-1\).  Above \(L+1\), the positive tail in
(TU.9) is Schwartz.  In the remaining strip, the kernel
and each fixed derivative are bounded by a fixed power of
\(1+L\).  Integrating its squared kernel therefore proves
\[
       \|UT_L\|_{\mathcal S_2(L^2,H^j)}
                    \leq C_j(1+L)^{d_j}\quad(j\geq0).
 \tag{TU.11}
\]
The same conclusion for the lower tail follows by writing
the canceled polynomial convolution first and estimating
the remaining Schwartz kernel; right multiplication by
the contraction \(W_L\) cannot increase that estimate.
Multiplication of distributions by the analytic \(L\)-factor
is legitimate and cancels every prescribed pole.  Thus the
closed multiplier realization gives the exact identity
\[
                         E_\chi T_L=C_aW_L .
 \tag{TU.12}
\]

We also need its ordinary Fourier transform in raw space.
On a fixed parity sector, the unitary raw Fourier transform
in Mellin coordinates is reflection \(\tau\mapsto-\tau\)
followed by a smooth unimodular multiplier on the real line.
For example these multipliers are
\[
 2(2\pi)^{-s}\Gamma(s)\cos(\pi s/2),\qquad
 -2i(2\pi)^{-s}\Gamma(s)\sin(\pi s/2),
                   \quad s=\tfrac12+i\tau ,
 \tag{TU.13}
\]
in the even and odd conventions, respectively.
The Mellin integral of cosine or sine first proves the
formula in its convergence strip; unitarity proves its
extension to the critical line.
Multiply \(b(-\tau)\) by this smooth multiplier.
Changing its finite pole prescriptions back to those used
in (TU.9) adds only delta derivatives at the reflected
zero ordinates.  These vanish on the reflected input
\(JW_L\), where \(Jf(t)=f(-t)\).
Consequently
\[
       \mathcal F_\epsilon T_L=U^{-1}C_{b^\#}JW_L
 \tag{TU.14}
\]
for a compactly supported distribution \(b^\#\) whose
kernel has precisely the two properties (TU.9).
This is an exact identity, not an approximation of a pole.

\emph{The two boundary families are compact in the required graph norms.}
Use the fixed smooth \(P=P_\Lambda^p\), \(Q=Q_R^p\) of
AC-T38.  The raw upper tail \((1-P)T_L\), after unitary
raw dilation by \(\Lambda\) and input translation by \(L\),
has, before the finite jet correction, the log kernel
\[
        (1-p(e^t))\,\kappa(t-s)\,
                              1_{[-2L,0]}(s).
 \tag{TU.15}
\]
Its output coordinate is bounded below by a fixed constant.
Since \(s\leq0\), (TU.9) makes this a Schwartz kernel in
the quadrant towards infinity.
It converges in Hilbert--Schmidt norm, with every fixed
output log derivative and every ordinary raw derivative,
to the same kernel with \(1_{(-\infty,0]}(s)\).
For ordinary raw derivatives use
\(\partial_u=u^{-1}\partial_{\log u}\), and recall
that the output raw coordinate is bounded away from zero.
All the resulting kernels are square integrable uniformly.

The finite jet correction tends to zero in these same norms.
In the translated input coordinate a normalized jet is
\(O(L^{-1/2})(1+|s|/L)^j\) on \([-2L,0]\), up to its
unit-modulus phase and uniformly bounded Gram coefficients.
The limiting kernel in (TU.15) has integrable moments of
every order in \(s\).
Cauchy--Schwarz, or direct integration of those moments,
therefore bounds each corrected column by \(O(L^{-1/2})\).
There are only finitely many columns.

The Fourier upper tail
\((1-p(\xi/R))\mathcal F_\epsilon T_L\) has the identical
description using (TU.14).
Its raw endpoint is \(\log R=L-\log q\), so only a fixed
translation by \(\log q\) occurs in its kernel.
It has the same Hilbert--Schmidt convergence and all the
same derivative bounds.

These facts also prove the two cross estimates
\[
 \|(1-Q)(1-P)T_L\|_{\mathcal S_2(H,H_D^j)}
 +\|(1-P)(1-Q)T_L\|_{\mathcal S_2(H,H_D^j)}
                   =O_{j,M}(c^{-M})\quad(j,M\geq0).
 \tag{TU.16}
\]
For the first estimate, take the Fourier tail of the
rescaled kernel (TU.15).  Every ordinary output derivative
is uniformly Hilbert--Schmidt, so its Fourier tail above
frequency \(c\) is \(O_M(c^{-M})\).
Applying a log derivative commutes with Fourier transform
by changing the sign of the dilation generator; derivatives
of the smooth cutoff are uniformly bounded in log coordinates.
The same reasoning applied to each derivative proves the
stated \(H_D^j\) bound.
For the second estimate use the Fourier-dual boundary family
and invert Fourier transform.  This proves (TU.16) directly
with the actual conductor scale.

\emph{Exact decomposition of the target residual.}
Let \(S_+=1_{(-\infty,L]}\) and \(S_-=1_{[-L,\infty)}\),
so \(S_L=S_+S_-\).
Raw time support gives \(E_\chi P\) target support at most
\(\Lambda\).  Raw Fourier support and primitive Poisson
give \(E_\chi Q\) target support at least
\(1/(qR)=1/\Lambda\).
All raw terms here belong to \(H_D^1\), by (TU.11) and
boundedness of the smooth log cutoffs on this space.
The support assertions first hold on the corresponding
time- or band-supported smooth cores.  Cutting off near
raw zero and then smoothing within the stated strict margin
approximates in \(H_D^{5/16}\); Fourier conjugation preserves
that norm.  Thus (TU.5) extends the support assertions to
the actual vectors in the following identity.
Consequently, with no commutation of \(P,Q\),
\[
 \begin{split}
 S_LE_\chi T_L-E_\chi PQPT_L
  ={}&S_LE_\chi(1-P)T_L\\
    &+S_LE_\chi(1-Q)PT_L
      +S_+E_\chi(1-P)QPT_L .
 \end{split}
 \tag{TU.17}
\]
For example, the last step of the expansion is
\((S_L-1)E_\chi QPT_L+E_\chi(1-P)QPT_L\);
its lower part is zero by the just stated Fourier gap,
and its upper part cancels because \(E_\chi P\) has
the upper support bound.  This verifies (TU.17).

Equation (TU.16), the boundedness of smooth log cutoffs
on \(H_D^1\), and (TU.5) permit replacing \(PT_L\) by
\(T_L\) in the second term, and replacing \(QPT_L\) by
\(PT_L\) in the third term, with Hilbert--Schmidt error
tending to zero faster than every power of \(c\).
For the latter replacement the exact raw error is
\[
 \begin{split}
 (1-P)(Q-1)PT_L
   ={}&-(1-P)(1-Q)T_L\\
      &+(1-P)(1-Q)(1-P)T_L.
 \end{split}
 \tag{TU.16a}
\]
The first term is the second cross estimate in (TU.16);
the second is the first cross estimate followed by the
uniformly \(H_D^1\)-bounded multiplier \(1-P\).
By (TU.15), the first and third terms, translated to the
upper target endpoint \(L\), converge in Hilbert--Schmidt
norm to fixed operators with range in \(L^2((-\infty,0])\).
The far lower target cutoff tends strongly to the identity
on their fixed Hilbert--Schmidt limits.
For the second term use the exact theta--Fourier identity
from AC-T25/AC-T32:
\(E_\chi\mathcal F_\epsilon^{-1}\) is a unit phase times
\(J_qE_{\bar\chi}\), with \(J_qF(x)=F(1/(qx))\).
The raw Fourier endpoint \(R\) then becomes the target
endpoint \(1/(qR)=1/\Lambda\).
Thus after translation to \(-L\) this second term converges
in Hilbert--Schmidt norm to a fixed operator with range
in \(L^2([0,\infty))\).
The graph-norm convergence of the raw boundary kernels
justifies applying these closed theta operators before
taking the limits.

\emph{Absorption of a fixed boundary operator by actual finite images.}
Under \(\mathsf A\), AC-T24, using Burnol's actual
nonprincipal sector with multiplier \(V=1\), says that
\[
 \overline{\{E_\chi\phi:
          \phi\text{ smooth of the required parity},
                    \operatorname{supp}\phi\Subset(-1,1)\}}
                  =L^2((-\infty,0])
 \tag{TU.18}
\]
in target log coordinates.
The same assertion holds for the dual character.
Smooth profiles may be chosen to vanish near raw zero:
cutting a smooth profile there changes its \(H_D^s\) norm
by a quantity tending to zero for every fixed \(s<1/2\),
as follows directly from its \(O(e^{t/2})\) log tail.
Equation (TU.5) then preserves density.

If a fixed Hilbert--Schmidt operator has range in the
right side of (TU.18), first approximate it by a finite-rank
operator, and then approximate each of its finitely many
output vectors by \(E_\chi\phi\) in (TU.18).
This approximates the operator in Hilbert--Schmidt norm.
Dilate those finitely many raw profiles to the time
endpoint \(\Lambda\).
AC-T38 shows that replacing each by its actual
\(\Pi_\Lambda\) projection changes its theta image by
a quantity tending to zero.
Therefore, for every upper boundary family just obtained,
\[
                 \|(1-Q_\Lambda)B_{\Lambda,+}\|_2\to0.
 \tag{TU.19}
\]
For a lower boundary family use the dual version of
(TU.18), Fourier transform its profiles back to the raw
space, and dilate their Fourier support to radius \(R\).
The second profile assertion of AC-T38 applies.
Primitive Poisson supplies exactly the lower target endpoint
\(1/(qR)\), so the same conclusion holds for that family.
All limits here keep the finite frequency interval and its
finite zero set fixed.  No completed span of all zero jets
in a fixed window is used.

By (TU.11) and AC-T38, the bulk \(E_\chi PQPT_L\) itself
is within \(o(1)\), in Hilbert--Schmidt norm, of actual
columns in \(E_\chi\Pi_\Lambda\).
Equations (TU.12), (TU.17)--(TU.19) now give
\[
                    \|(1-Q_\Lambda)S_LC_aW_L\|_2\to0.
 \tag{TU.20}
\]
There is one additional input boundary:
\(S_LC_a(1-S_L)\).
Its convolution kernel is Schwartz because \(a\) is
smooth compactly supported.
Splitting its input above \(L\) and below \(-L\),
and translating each endpoint to zero, gives two fixed
quadrant Hilbert--Schmidt limits, with outputs on the
appropriate target half-lines.
The same boundary absorption argument proves
\[
                  \|(1-Q_\Lambda)S_LC_a(1-S_L)\|_2\to0.
 \tag{TU.21}
\]

\emph{The matching upper bound.}
The positive defect has the exact distance formula
\[
       \operatorname{Tr}(D_\Lambda C_aC_a^*)
                    =\|(1-Q_\Lambda)S_LC_a\|_2^2 .
 \tag{TU.22}
\]
Decompose its input orthogonally into \(Z_L\) and its
orthogonal complement.
Equations (TU.20)--(TU.21) show that the second summand
tends to zero.  The first is bounded above by
\(\|S_LC_aP_{Z,L}\|_2^2\), whose limit is (TU.8).
This proves the required upper limit.
The lower limit is the finite-jet argument of AC-T31:
that proof uses only boundedness and continuity of the
multiplier at the finitely many ordinates, so it applies
to the present \(a\) verbatim.
This proves (TU.6).

Finally polarization gives convergence against every
smooth compact frequency test: write such a test as
\(a\overline b\), choosing \(a=1\) on its support.
For each compact set choose a smooth \(a\) equal to one
there.  Equation (TU.6) bounds the local masses uniformly
for large \(\Lambda\).
Approximation of continuous compact tests by smooth ones
then gives the stated vague convergence.
\end{proof}

\begin{remark}[What AC-T39 leaves in the original trace problem]
The local upper bound rules out an excess positive measure
on every fixed compact Mellin-frequency interval.
The original cutoff trace tests have rapidly decreasing,
but generally noncompact, Mellin transforms.
Their full limit still requires a uniform weighted tail
estimate for \(\mu_\Lambda\), or an equivalent
justified passage from these compact frequency tests.
Local convergence alone does not supply that passage.
AC-T40--T41 subsequently supply this estimate and complete
the full trace in the fixed nonprincipal primitive
Dirichlet prolate sectors considered here.
The exact actual-adelic cutoff dictionary and the other
scope qualifications of AC-T11 and AC-T30 remain in force.
All approximations in this section concern actual
Hilbert--Schmidt column operators and orthogonal projections.
Burnol's Theorem 3.2 is used only through the already proved
half-line closure (TU.18), with its nonprincipal multiplier
identically one.  The finite-pole inverse, strict cutoff
margin, and boundary approximation supply the further
finite-prolate comparison here.
No lifting of a Gram inverse to an LF representation,
Grothendieck nuclear trace identification, or new adelic
cutoff equivalence is asserted by this argument.
\end{remark}
